\documentclass[10pt,a4paper]{amsart}
\usepackage{amsmath,amssymb,amsthm,graphicx}
% Retained minimal preamble for the Stacks Project chapter more-algebra.tex (More on Algebra).
% The chapter file pinned-original-source.tex carries only \input{preamble} on its first line, so no
% preamble text lives inside the chapter itself.  The Stacks Project uses one single shared preamble for
% every chapter; the snapshot stages it identically as sources/stacks-chow/preamble.tex and
% sources/stacks-etale-cohomology/preamble.tex.  The block below is copied mechanically, line for line,
% from sources/stacks-chow/preamble.tex lines 14-18 (the Xy-pic setup) and 154-217 (the native theorem
% declarations together with the mathematical macros).  Nothing was added, removed or redefined; the
% wrapper class and the amsmath/amssymb/amsthm/graphicx base packages are supplied by the assembler.
% Chapter: more-algebra.tex -- More on Algebra.
% For commutative diagrams we use Xy-pic
\usepackage[all]{xy}

% We use 2cell for 2-commutative diagrams.
\xyoption{2cell}

% Theorem environments.
%
\theoremstyle{plain}
\newtheorem{theorem}[subsection]{Theorem}
\newtheorem{proposition}[subsection]{Proposition}
\newtheorem{lemma}[subsection]{Lemma}

\theoremstyle{definition}
\newtheorem{definition}[subsection]{Definition}
\newtheorem{example}[subsection]{Example}
\newtheorem{exercise}[subsection]{Exercise}
\newtheorem{situation}[subsection]{Situation}

\theoremstyle{remark}
\newtheorem{remark}[subsection]{Remark}
\newtheorem{remarks}[subsection]{Remarks}

\numberwithin{equation}{subsection}

% Macros
%
\def\lim{\mathop{\mathrm{lim}}\nolimits}
\def\colim{\mathop{\mathrm{colim}}\nolimits}
\def\Spec{\mathop{\mathrm{Spec}}}
\def\Hom{\mathop{\mathrm{Hom}}\nolimits}
\def\Ext{\mathop{\mathrm{Ext}}\nolimits}
\def\SheafHom{\mathop{\mathcal{H}\!\mathit{om}}\nolimits}
\def\SheafExt{\mathop{\mathcal{E}\!\mathit{xt}}\nolimits}
\def\Sch{\mathit{Sch}}
\def\Mor{\mathop{\mathrm{Mor}}\nolimits}
\def\Ob{\mathop{\mathrm{Ob}}\nolimits}
\def\Sh{\mathop{\mathit{Sh}}\nolimits}
\def\NL{\mathop{N\!L}\nolimits}
\def\CH{\mathop{\mathrm{CH}}\nolimits}
\def\proetale{{pro\text{-}\acute{e}tale}}
\def\etale{{\acute{e}tale}}
\def\QCoh{\mathit{QCoh}}
\def\Ker{\mathop{\mathrm{Ker}}}
\def\Im{\mathop{\mathrm{Im}}}
\def\Coker{\mathop{\mathrm{Coker}}}
\def\Coim{\mathop{\mathrm{Coim}}}

% Boxtimes
%
\DeclareMathSymbol{\boxtimes}{\mathbin}{AMSa}{"02}

%
% Macros for moduli stacks/spaces
%
\def\QCohstack{\mathcal{QC}\!\mathit{oh}}
\def\Cohstack{\mathcal{C}\!\mathit{oh}}
\def\Spacesstack{\mathcal{S}\!\mathit{paces}}
\def\Quotfunctor{\mathrm{Quot}}
\def\Hilbfunctor{\mathrm{Hilb}}
\def\Curvesstack{\mathcal{C}\!\mathit{urves}}
\def\Polarizedstack{\mathcal{P}\!\mathit{olarized}}
\def\Complexesstack{\mathcal{C}\!\mathit{omplexes}}
% \Pic is the operator that assigns to X its picard group, usage \Pic(X)
% \Picardstack_{X/B} denotes the Picard stack of X over B
% \Picardfunctor_{X/B} denotes the Picard functor of X over B
\def\Pic{\mathop{\mathrm{Pic}}\nolimits}
\def\Picardstack{\mathcal{P}\!\mathit{ic}}
\def\Picardfunctor{\mathrm{Pic}}
\def\Deformationcategory{\mathcal{D}\!\mathit{ef}}

% --- appended 2026-10-12 (Chief): the project own definitions of the reference/slogan/history
% display environments, copied mechanically from the shared Stacks preamble lines 8-12
% (\usepackage{verbatim} plus the three \newenvironment lines, which the chapter preamble keeps
% verbatim from the project preamble). Nothing added, removed or redefined: these environments
% carry no mathematical content and the project itself typesets them as comments. Without these
% three lines a counted statement whose own text uses them fails with Environment slogan undefined.
\usepackage{verbatim}
\newenvironment{reference}{\comment}{\endcomment}
\newenvironment{slogan}{\comment}{\endcomment}
\newenvironment{history}{\comment}{\endcomment}
\usepackage{cleveref}
\crefname{originallemma}{original lemma}{original lemmas}
\Crefname{originallemma}{Original lemma}{Original lemmas}
\makeatletter
\newlabel{algebra-lemma-cover}{{00EO}{source}{}{sourcepointer-1}{}}
\newlabel{algebra-lemma-cover@cref}{{[originallemma][0][]00EO}{[1][0][]source}{}{}{}}
\makeatother
\crefname{originallemma}{original lemma}{original lemmas}
\Crefname{originallemma}{Original lemma}{Original lemmas}
\makeatletter
\newlabel{algebra-lemma-extension}{{0519}{source}{}{sourcepointer-2}{}}
\newlabel{algebra-lemma-extension@cref}{{[originallemma][0][]0519}{[1][0][]source}{}{}{}}
\makeatother
\crefname{originallemma}{original lemma}{original lemmas}
\Crefname{originallemma}{Original lemma}{Original lemmas}
\makeatletter
\newlabel{lemma-localize-relative-finite-presentation}{{21015}{source}{}{sourcepointer-3}{}}
\newlabel{lemma-localize-relative-finite-presentation@cref}{{[originallemma][0][]21015}{[1][0][]source}{}{}{}}
\makeatother
\crefname{originallemma}{original lemma}{original lemmas}
\Crefname{originallemma}{Original lemma}{Original lemmas}
\makeatletter
\newlabel{lemma-relatively-finitely-presented}{{20923}{source}{}{sourcepointer-4}{}}
\newlabel{lemma-relatively-finitely-presented@cref}{{[originallemma][0][]20923}{[1][0][]source}{}{}{}}
\makeatother
\title{Finite presentation relative to a base ring can be glued along a finite cover by principal opens}
\author{The Stacks Project Authors}
\date{stacks-project:more-algebra.tex:local-snapshot-20261012; GNU FDL 1.2 or later}
\begin{document}
\maketitle
\setcounter{subsection}{1}
\noindent\textit{Editorial note.} All delivered mathematical text is the author's own, copied from the pinned Stacks Project chapter file more-algebra.tex; no statement or proof step was rewritten, reordered or removed.\par\medskip
\noindent\textit{Editorial note.} The chapter file pinned-original-source.tex carries only \textbackslash{}input\{preamble\} on its first line, so no preamble text lives inside the chapter. The retained file sources/stacks-more-algebra/retained-preamble.tex is a mechanical, line-for-line copy of the project's single shared preamble as staged in the snapshot; nothing was added, removed or redefined and no package was added.\par\medskip
\noindent\textit{Editorial note.} The source URL is the source-level Stacks Project URL, as required when more than one card is produced from one source; the per-claim tag is carried by claim\_key.\par\medskip
\noindent\textit{Editorial note.} The author writes a period before the sentence opening 'such that y\_i maps to f\_i' at proof line 21146; the punctuation is reproduced unchanged. The five references inside the counted proof span are recorded as non-load-bearing navigation pointers.\par\medskip
\noindent\textit{Allowed standard prerequisites (editorial).}\begin{itemize}
\item Finite presentation of a module relative to a base ring and its behaviour under localization. used at proof lines 21142-21143 and 21154-21157. Is the property being glued and the input that makes the localized modules finitely presented over the localized rings.
\item The cover criterion for finite presentation of a module. invoked at proof lines 21157-21158. Turns the local finite presentation of the M\_\{f\_i\} into finite presentation of M.
\item The auxiliary algebra P = R[x, y, z]/(sum y\_i z\_i - 1) realizing the cover generators and their unit relation. constructed at proof lines 21145-21153. Provides the ring over which M becomes a finite module and in which the final vanishing argument is run.
\item A module with vanishing localizations at elements generating the unit ideal is zero. used at proof lines 21171-21174. Is the final step closing K/K' = 0.
\end{itemize}\medskip
\ifdefined\hypertarget\hypertarget{sourcepointer-1}{}\fi
\noindent\textit{Source reference.} Original lemma 00EO: focused-seeds/AGP/sources/stacks-snapshot-20261011/tags.txt, source lines 542--542. Cross-chapter navigation pointer into the Stacks Project; the standard fact it supplies is applied inside the counted proof span and the argument resting on it stays there.\par\medskip
\ifdefined\hypertarget\hypertarget{sourcepointer-2}{}\fi
\noindent\textit{Source reference.} Original lemma 0519: focused-seeds/AGP/sources/stacks-snapshot-20261011/tags.txt, source lines 6527--6527. Cross-chapter navigation pointer into the Stacks Project; the standard fact it supplies is applied inside the counted proof span and the argument resting on it stays there.\par\medskip
\ifdefined\hypertarget\hypertarget{sourcepointer-3}{}\fi
\noindent\textit{Source reference.} Original lemma 21015: sources/stacks-more-algebra/pinned-original-source.tex, source lines 21015--21015. Same-chapter navigation pointer; the statement it names is standard for the section and the argument that consumes it stays inside the counted proof span.\par\medskip
\ifdefined\hypertarget\hypertarget{sourcepointer-4}{}\fi
\noindent\textit{Source reference.} Original lemma 20923: sources/stacks-more-algebra/pinned-original-source.tex, source lines 20923--20923. Same-chapter navigation pointer; the statement it names is standard for the section and the argument that consumes it stays inside the counted proof span.\par\medskip
\begin{lemma}
\label{lemma-glue-relative-finite-presentation}
Let $R \to A$ be a finite type ring map. Let $M$ be an $A$-module.
Let $f_1, \ldots, f_r \in A$ generate the unit ideal.
The following are equivalent
\begin{enumerate}
\item each $M_{f_i}$ is finitely presented relative to $R$, and
\item $M$ is finitely presented relative to $R$.
\end{enumerate}
\end{lemma}

\begin{proof}
The implication (2) $\Rightarrow$ (1) is in
Lemma \ref{lemma-localize-relative-finite-presentation}.
Assume (1). Write $1 = \sum f_ig_i$ in $A$.
Choose a surjection
$R[x_1, \ldots, x_n, y_1, \ldots, y_r, z_1, \ldots, z_r] \to A$.
such that $y_i$ maps to $f_i$ and $z_i$ maps to $g_i$. Then we
see that there exists a surjection
$$
P = R[x_1, \ldots, x_n, y_1, \ldots, y_r, z_1, \ldots, z_r]/(\sum y_iz_i - 1)
\longrightarrow
A.
$$
By
Lemma \ref{lemma-relatively-finitely-presented}
we see that $M_{f_i}$ is a finitely presented $A_{f_i}$-module, hence by
Algebra, Lemma \ref{algebra-lemma-cover}
we see that $M$ is a finitely presented $A$-module.
Hence $M$ is a finite $P$-module (with $P$ as above).
Choose a surjection $P^{\oplus t} \to M$.
We have to show that the kernel $K$ of this map is a finite
$P$-module. Since $P_{y_i}$ surjects onto
$A_{f_i}$ we see by
Lemma \ref{lemma-relatively-finitely-presented}
and
Algebra, Lemma \ref{algebra-lemma-extension}
that the localization $K_{y_i}$ is a finitely generated
$P_{y_i}$-module. Choose elements
$k_{i, j} \in K$, $i = 1, \ldots, r$, $j = 1, \ldots, s_i$ such
that the images of $k_{i, j}$ in $K_{y_i}$ generate.
Set $K' \subset K$ equal to the $P$-module
generated by the elements $k_{i, j}$. Then $K/K'$ is a module
whose localization at $y_i$ is zero for all $i$. Since $(y_1, \ldots, y_r) = P$
we see that $K/K' = 0$ as desired.
\end{proof}

\end{document}
